%% nilHMM methods paper — mathematical supplement.
%% Comprehensive, connected derivation from the low-coverage foundations through
%% the shared HMM engine to simcross calibration and GWAS-landscape geometry.
%% Structure modeled on the RTIGER supplement (Campos-Martin et al. 2023):
%% numbered Texts, full derivations, one narrative. Compiles standalone or is
%% \input into main.tex's \section*{Supplement}.
\documentclass[11pt]{article}
\usepackage{amsmath,amssymb}
\usepackage{graphicx}
\usepackage[margin=1in]{geometry}
\usepackage{booktabs}
\usepackage{array}
\usepackage[round,authoryear]{natbib}
\usepackage[colorlinks=true,linkcolor=blue,urlcolor=blue,citecolor=blue]{hyperref}
\setlength{\emergencystretch}{3em}

% ---- macros ----
\DeclareMathOperator*{\argmax}{argmax}
\DeclareMathOperator*{\argmin}{argmin}
\newcommand{\REF}{\mathrm{REF}}
\newcommand{\HET}{\mathrm{HET}}
\newcommand{\ALT}{\mathrm{ALT}}
\newcommand{\Oo}{\mathcal{O}}
\newcommand{\Ss}{\mathcal{S}}
\newcommand{\Sr}{S_r}
\newcommand{\Pois}{\mathrm{Pois}}
\newcommand{\BetaBin}{\mathrm{BetaBin}}
\newcommand{\Mult}{\mathrm{Multinom}}
\newcommand{\Bfun}{\mathrm{B}}                     % beta function
\newcommand{\miss}{\mathrm{miss}}
\newcommand{\cov}{\lambda}                          % per-marker coverage
\newcommand{\code}[1]{\texttt{#1}}

\title{Supplemental Texts\\[2pt]
\large Simulation-calibrated HMM ancestry calling for near-isogenic lines\\
across sequencing modalities (nilHMM)}
\author{Rodr\'iguez-Zapata \textit{et al.}}
\date{}

\begin{document}
\maketitle

\medskip
\begin{description}
\item[Text S1] Foundations --- the low-coverage regime and its missing-data model
\item[Text S2] Foundations --- genotype calling under low coverage
\item[Text S3] The unified duration-aware three-state HMM engine
\item[Text S4] The caller family
\item[Text S5] The \code{simcross} model and single-locus validation
\item[Text S6] Two-stage simulation calibration
\item[Text S7] GWAS-landscape geometry
\item[Text S8] Rigid-HMM derivations (the RTIGER port)
\item[Text S9] Pedigree-coupled ancestry calling
\end{description}
\clearpage

%======================================================================
\section{Foundations --- the low-coverage regime and its missing-data model}
\label{sec:foundations-missing}
%======================================================================
% Motivation: the ZEAL/BZea population is genuinely low-coverage skim data; what
% makes it tractable is that per-marker missingness obeys a clean
% exponential-with-floor law in coverage. This section builds that law up from a
% Poisson baseline, fits it to BZea BAMs, shows it holds across modalities, and
% establishes the scale-invariance that lets us reuse it as the simcross
% degradation model.

The BZea/ZEAL population is genotyped by genome-skim sequencing: roughly
$1{,}400$ near-isogenic lines sequenced to a nominal $0.8\times$ against the
$2.2$\,Gb B73 NAM v5 reference with $150$\,bp reads (about $1.2\times10^{7}$
reads per sample). At this depth most sites in most samples carry no read
\citep{elshire_2011,nielsen_2011}, so whether ancestry can be called at all rests
on \emph{how} missingness scales with coverage. This section derives that scaling law, fits it to the real BZea
alignments, shows it is stable across sequencing modalities, and establishes the
scale-invariance that later lets us reuse it as the \texttt{simcross} degradation
model (Text~S5--S6).

\subsection{The Poisson coverage baseline}
\label{sec:poisson}

If reads land independently and uniformly, the number of reads $x$ covering a
given site is Poisson with mean $\cov$, the sequencing redundancy (mean coverage)
--- the Lander--Waterman model of random read placement \citep{lander_waterman_1988,sims_2014} ---
\begin{equation}
P(x) = \frac{\cov^{x} e^{-\cov}}{x!},
\qquad \cov = \frac{L N}{G},
\label{eq:poisson}
\end{equation}
where $G$ is the genome length in bp, $L$ the read length, and $N$ the number of
reads. A site is missing when $x=0$, so the baseline per-site missingness and the
covered fraction are
\begin{equation}
P_{\miss}^{\mathrm{Pois}}(\cov) = e^{-\cov},
\qquad
P_{\mathrm{cov}}(\cov) = 1 - e^{-\cov}.
\label{eq:poisson-missing}
\end{equation}
Equation~\eqref{eq:poisson-missing} is the random-sampling ideal: missingness
decays to zero as coverage grows, limited only by shot noise. Real skim data does
not behave this way --- a fraction of the genome stays uncallable at \emph{any}
depth --- which motivates an irreducible floor.

\subsection{Exponential-with-floor missingness, fit to BZea}
\label{sec:floor}

We model per-site missingness as a mixture of a coverage-driven Poisson term and a
coverage-independent floor (a ``structural floor plus sampling'' two-component form
shared with ecological occupancy--detection models, \citealp{mackenzie_2002}),
\begin{equation}
\boxed{\;\miss(\cov) \;=\; \pi \;+\; (1-\pi)\,e^{-k\cov}\;}
\label{eq:floor}
\end{equation}
with two parameters: the \emph{floor} $\pi\in[0,1]$, the fraction of sites that
stay uncalled even at high coverage (repetitive DNA, structural variants, and
highly polymorphic regions where short reads fail to map or are filtered ---
low mappability / the inaccessible genome, \citealp{li_maq_2008,derrien_2012,tgp_2015}),
and the
\emph{decay rate} $k>0$, which sets how fast the callable fraction fills in with
coverage relative to the Poisson ideal ($k=1$ recovers
Eq.~\eqref{eq:poisson-missing} up to the floor; the product $k\cov$ acts as an
effective coverage over informative variants).

We fit Eq.~\eqref{eq:floor} directly to the BZea \emph{Wideseq} BAM alignments
\emph{before} genotype calling, to avoid confounding the missingness law with
downstream VCF filters. Per-site, per-sample allele read counts
(GATK \texttt{CollectAllelicCounts}) were taken at the ${\sim}27$ million known
teosinte-vs-B73 variant sites and aggregated into $1$\,Mb bins; per sample,
$\cov = \texttt{DEPTH\_SUM}/\texttt{VARIANT\_COUNT}$ and
$\miss = 1 - \texttt{INFORMATIVE\_VARIANT\_COUNT}/\texttt{VARIANT\_COUNT}$. The
realized mean coverage is $\cov\approx0.59$ (below the $0.8\times$ design goal),
at which the observed missingness is $\approx0.67$ whereas the Poisson baseline
predicts only $e^{-0.59}\approx0.554$ --- Poisson systematically under-predicts,
as expected when read depth is over-dispersed by PCR and GC-content bias
\citep{benjamini_speed_2012,ross_2013}.
A nonlinear least-squares fit of Eq.~\eqref{eq:floor}
(\texttt{nlsLM}, $\pi\in[0,1]$, $k\in[0,10]$) gives
\begin{equation}
\pi \approx 0.346,
\qquad
k \approx 1.256
\label{eq:floor-fit-bzea}
\end{equation}
i.e.\ about one-third of sites are unmappable at any depth --- a floor expected to
be large in maize, whose genome is ${\sim}85\%$ transposable-element and
structurally variable sequence \citep{schnable_2009,gui_2022}. Fit quality is judged
by concordance rather than a summary statistic: on an
observed-vs-Poisson-expected plot the points lie systematically above the $1{:}1$
line, whereas against the floor-model expectation they collapse onto it.
\begin{figure}[t]
\centering
% TODO(export from missing-data-floor-model.qmd: final_plot_poisson + final_plot_floor)
\caption{\textbf{Missingness in BZea Wideseq follows an exponential-with-floor
law.} Observed per-bin missingness against the Poisson expectation
$e^{-\cov}$ (left; points fall above the $1{:}1$ line --- Poisson under-predicts)
and against the floor-model expectation $\pi+(1-\pi)e^{-k\cov}$ (right; points on
the $1{:}1$ line), colored by coverage. Dashed vertical line: the fitted floor
$\pi\approx0.346$.}
\label{fig:floor-fit}
\end{figure}

\subsection{Missingness across modalities}
\label{sec:model-comparison}

The same law, fit to the same BZea samples viewed as three data products ---
whole-genome (WGS, \texttt{picard CollectWgsMetrics}), Wideseq (the $27$M
teosinte-variant panel), and the curated SNP50K panel --- gives a floor that
falls monotonically with SNP ascertainment while the decay rate relaxes toward the
Poisson value $k\to1$ as the panel is cleaned (Table~\ref{tab:modality}). That
$\pi$ tracks mappability exactly as its mechanistic reading predicts --- filtering
to informative, well-mapped sites removes irreducible missingness --- is evidence
the floor model captures a real property of the data rather than a curve-fitting
artifact.
\begin{table}[t]
\centering
\begin{tabular}{l c}
\toprule
dataset (same BZea samples) & floor $\pi$ \\
\midrule
WGS --- whole genome (unfiltered)      & $0.446$ \\
Wideseq --- $27$M teosinte variants    & $0.346$ \\
SNP50K --- curated $50$K panel         & $0.161$ \\
\bottomrule
\end{tabular}
\caption{Fitted floor $\pi$ across sequencing modalities. The floor decreases with
ascertainment (WGS $>$ Wideseq $>$ SNP50K); the decay $k$ likewise relaxes toward
the Poisson value $1$ as the panel is cleaned.}
\label{tab:modality}
\end{table}

\subsection{Scale-invariance of $\miss(\cov)$}
\label{sec:ergodicity}

The degradation benchmark (Text~S5) applies the missingness law \emph{per $1$\,Mb
bin}, but Eq.~\eqref{eq:floor-fit-bzea} was fit \emph{between samples}. This is
licensed only if $\miss(\cov)$ is the same function whether $\cov$ varies because
one picked a low-depth \emph{sample} or a low-mappability \emph{region} --- i.e.\
if $\cov$ is a near-sufficient statistic for missingness. Fitting
Eq.~\eqref{eq:floor} at both scales on real data confirms this
(Table~\ref{tab:ergodicity}): the within-genome (per-bin) and between-sample
curves nearly coincide ($\Delta\pi=0.033$, $\Delta k=0.080$), separating only in
the high-coverage tail ($\cov>1\times$). An ergodic-by-construction simulation
recovers the input parameters at both scales ($\Delta\pi\approx0$,
$\Delta k\approx0.001$), confirming the machinery; the small real residual (a
per-bin floor $0.313$ below the per-sample $0.346$) is the genuine spatial
structure --- pericentromeric and repeat-rich bins are systematically
low-mappability, so the genomic axis is not stationary and the transfer is
approximate, not exact. Within a bin the ${\sim}10^{3}$--$10^{4}$ informative sites
are exchangeable, and by the central limit theorem the bin statistic is governed
by the mean per-site rate ($\approx\cov$) rather than the per-site coverage shape,
which is what protects the transfer.
\begin{table}[t]
\centering
\begin{tabular}{@{}l l r c c@{}}
\toprule
data & scale & $n$ & $\pi$ & $k$ \\
\midrule
real       & within-genome ($1$\,Mb bins) & $3{,}057{,}288$ & $0.313$ & $1.176$ \\
real       & between-genome (samples)     & $1{,}434$       & $0.346$ & $1.256$ \\
simulation & within-genome (bins)         & $212{,}900$     & $0.346$ & $1.255$ \\
simulation & between-genome (samples)     & $100$           & $0.346$ & $1.256$ \\
\bottomrule
\end{tabular}
\caption{Scale-invariance of the missingness law. Fitting
Eq.~\eqref{eq:floor} within-genome (across $1$\,Mb bins) and between-genome
(across samples) gives nearly the same $(\pi,k)$ on real data; the simulation is a
control, ergodic by construction.}
\label{tab:ergodicity}
\end{table}

The practical consequence is the transfer function used throughout the rest of the
paper: because a single two-parameter law summarizes missingness as a function of
coverage, and that law is stable across data products and scales, a clean genome
can be \emph{analytically} degraded to any target coverage --- drawing a per-sample
depth $\cov_s\sim\mathrm{Normal}(0.590,\,0.226)$ (the between-sample coverage
distribution, favored over Gamma and lognormal by AIC) and applying the fitted
$(\pi,k)$ per bin --- rather than resampling real reads. The floor model is that
transfer function; the scale-invariance result is the license to apply it per bin
(Text~S5--S6).

%======================================================================
\section{Foundations --- genotype calling under low coverage}
\label{sec:foundations-calling}
%======================================================================
% Given the reads that survive, how is a call made, and why is low coverage hard?

Missingness (Text~S1) decides \emph{whether} a site is observed; this section
concerns the site once it is. The upshot is that at skim depth a single-marker
genotype call carries almost no zygosity information, and worse, that the
information ceiling is not lifted by adding reads or by a better prior --- which is
precisely why ancestry must be inferred as a spatially-correlated latent state
(Text~S3) rather than read off site by site.

\subsection{Single-read genotype likelihoods}
\label{sec:gl}

At a biallelic site let the reference (B73) allele be $A$ ($0$) and the alternate
(donor) allele $B$ ($1$), with genotypes $\{AA, AB, BB\}=\{0/0,0/1,1/1\}$ and a
per-base error $\varepsilon = 10^{-Q/10}$. The genotype likelihood of the reads
averages the two alleles of the diploid genotype with equal weight (the standard
genotype-likelihood model of NGS SNP callers, \citealp{depristo_2011,nielsen_2011}),
\begin{equation}
P(\text{reads}\mid G=\{A_1,A_2\})
= \prod_{i}\Big[\tfrac12 P(b_i\mid A_1) + \tfrac12 P(b_i\mid A_2)\Big],
\qquad
P(b\mid A)=
\begin{cases}
1-\varepsilon & b=A\\
\varepsilon/3 & b\ne A .
\end{cases}
\label{eq:gl}
\end{equation}
The equal-weight averaging pins the heterozygote likelihood of any single read
near $\tfrac12$, so one read carries at most a ${\sim}2{:}1$ likelihood ratio.
For a single ALT read (at $\varepsilon=0.01$):
\begin{center}
\begin{tabular}{l c c c c}
\toprule
$G$ & $P(\text{read}\mid G)$ & value & Phred $-10\log_{10}L$ & normalized PL \\
\midrule
$AA$ & $\varepsilon/3$                               & $0.0033$ & $24.8$ & $25$ \\
$AB$ & $\tfrac12\tfrac{\varepsilon}{3}+\tfrac12(1-\varepsilon)$ & $0.498$  & $3.0$  & $3$ \\
$BB$ & $1-\varepsilon$                               & $0.99$   & $0.04$ & $0$ \\
\bottomrule
\end{tabular}
\end{center}
The genotype quality $\mathrm{GQ}=(\text{2nd-smallest PL})-(\text{smallest PL})=3$:
one ALT read is a $2{:}1$ ($\approx 3$-PL) preference for donor-hom over het --- the
entire information content of one read, identical across callers because
Eq.~\eqref{eq:gl} is.

\subsection{The prior decides the call, and accumulation only helps a shared state}
\label{sec:prior}

Every posterior caller multiplies the (caller-identical) likelihood by a
genotype prior and reports the MAP genotype \citep{nielsen_2011,korneliussen_angsd_2014},
\begin{equation}
P(G\mid\text{reads}) \;\propto\;
\underbrace{P(\text{reads}\mid G)}_{\text{GL --- identical across tools}}
\;\times\;
\underbrace{\pi(G)}_{\text{prior --- swappable}},
\qquad
\hat G=\argmax_G P(\text{reads}\mid G)\,\pi(G).
\label{eq:posterior}
\end{equation}
Under a Hardy--Weinberg prior $\pi=\big((1-p)^2,\,2p(1-p),\,p^2\big)$ at the
cohort's mean alternate-allele frequency $\bar p=0.0152$, the single-ALT-read
posterior favors het over donor-hom by
\begin{equation}
\frac{P(AB\mid\text{read})}{P(BB\mid\text{read})}
=\frac{p(1-p)}{(1-\varepsilon)\,p^{2}}
=\frac{1-p}{(1-\varepsilon)\,p}\approx\frac{1}{p}\approx 66 .
\label{eq:hwe-ratio}
\end{equation}
Because every panel site is rare ($p\ll\tfrac12$; no site has $p\ge0.1$), a lone
ALT read is forced to het by a factor ${\sim}66$ --- the mechanical origin of the
observed skim genotype matrix ($97\%$ REF-hom, $\sim3\%$ het, essentially $0\%$
donor-hom; $448{,}805$ het vs.\ $355$ donor-hom calls at depth~1). Swapping the
prior to the breeding-design frequencies (BC2S3: $f_{AA}=0.859$, $f_{AB}=0.0312$,
$f_{BB}=0.1094$; the donor allele enters at frequency $0.125$) flips the \emph{same}
read from het ($P=0.81$ under HWE) to donor-hom ($P=0.855$, $\approx 7{:}1$). The
prior only relabels the tie: the data still speaks at $\approx 3$~PL, so \emph{one
read cannot resolve het from hom under any prior}, and because
Eq.~\eqref{eq:hwe-ratio} depends on $p$, the bias persists as depth grows until $p$
approaches $\tfrac12$ --- it is not removed by adding reads.

What accumulation \emph{can} help is a \emph{shared} latent state $S$ (the ancestry
of a chromosomal stretch), whose posterior multiplies per-site emissions
\begin{equation}
P(S\mid\text{reads}_{1:n}) \;\propto\; \pi(S)\prod_{j=1}^{n}
\underbrace{\Big[\textstyle\sum_{G} P(\text{reads}_j\mid G)\,P(G\mid S)\Big]}_{\text{emission}_j},
\label{eq:accumulate}
\end{equation}
and for a biallelic site the count pair $(n_{\mathrm{ref}},n_{\mathrm{alt}})$ is a
sufficient statistic of the per-site genotype likelihood --- so the count
(BetaBinomial) emission of Text~S3 is a direct parameterization of
Eq.~\eqref{eq:accumulate}, with overdispersion added for mapping/paralog noise.

\subsection{Dosage non-identifiability of the SNP50K skim}
\label{sec:identifiability}

Two independent mechanisms make per-site zygosity unrecoverable at skim depth;
neither is fixed by more reads or a better caller.

\paragraph{Rare-allele prior.} As in Eq.~\eqref{eq:hwe-ratio}, with the cohort
allele-frequency spectrum sitting entirely in the het branch
($\bar p=0.0152$; $41.8\%$ of sites at $p<0.01$, none at $p\ge0.1$), the het:donor-hom
posterior ratio $\approx 1/p\approx 66$ forces het regardless of which prior is
used, and persists with depth.

\paragraph{Outbred-donor dilution.} Teosinte donors are outbred and heterozygous,
and the panel sites are \emph{common species variants}, not accession-fixed alleles.
The gamete transmitted to the cross carries the donor ALT allele only with
probability equal to the donor's ALT-allele frequency at panel sites,
\begin{equation}
d \;=\; 0\cdot f_0 + \tfrac12 f_1 + 1\cdot f_2 \;=\; f_2 + \tfrac12 f_1
\;\approx\; 0.28\text{--}0.35,
\label{eq:donor-dose}
\end{equation}
with $f_0\approx60\%$ of sites not segregating in a given donor (so the gamete looks
like B73) and only about one site in three informative. This compresses the
state-conditional expected ALT fraction from the textbook $\{0,\tfrac12,1\}$ to
\begin{equation}
\mathbb{E}[\text{alt frac}\mid\text{state}] \approx
\{\,\varepsilon,\; \tfrac{d}{2},\; d\,\}
\approx\{\,0,\; 0.15,\; 0.30\,\},
\label{eq:compressed-means}
\end{equation}
so het and donor-hom sit at $\approx 0.15$ and $0.30$ --- a separation of only
$\approx 0.15$, barely resolvable. Reads confirm the compression: $2\times$(het
alt-fraction) should equal the donor-hom alt-fraction, but empirically donor-hom
exceeds it (e.g.\ \emph{Z.\ mexicana} $0.49$ vs.\ $0.37$), because at $0.4\times$ the
caller thresholds alt-fraction instead of resolving zygosity. Restoring
$\{0,\tfrac12,1\}$ would require restricting to sites fixed for the ALT allele
\emph{in that donor accession} ($f_2\to1$, $d\to1$), which needs the donor's own
genotypes --- unavailable. The usable levers are therefore all upstream (donor-
informative site selection, per-taxon fitted emission means, and the breeding
prior in place of HWE), not in the emission model itself.

\subsection{Emission by depth: counts vs.\ a hard call}
\label{sec:emission-by-depth}

Given that per-site zygosity is a \emph{data} limit, the remaining choice --- which
emission family to attach to the fixed three-state model --- is dictated by depth,
not preference (Table~\ref{tab:depth-regime}). The count (BetaBinomial) emission
earns its keep only in the intermediate band: below it the signal is degenerate and
no emission recovers information the reads do not contain; above it the count
likelihood collapses to a near-delta and a categorical hard-genotype call carries
the same information at lower cost (the BetaBinomial cost scales with the number of
distinct $(n,k)$ count pairs, i.e.\ with depth, and is capped around $20\times$).
Only the emission $p(\text{obs}\mid\text{state})$ swaps across regimes; the states
and the duration layer (Text~S3) are unchanged.
\begin{table}[t]
\centering
\small
\begin{tabular}{@{}l l l@{}}
\toprule
depth & regime & emission \\
\midrule
$\ll 1\times$ (skim, ${\sim}0.4\times$) & het invisible; data-limited & BetaBinomial counts (best available) \\
${\sim}1$--$20\times$ (BRB-seq, ${\sim}2.8\times$) & depth-weighting informative & \textbf{BetaBinomial counts} \\
$\gtrsim 20\times$ (target-seq / 45K) & saturated; likelihood $\to$ delta & categorical hard GT (equivalent, cheaper) \\
\bottomrule
\end{tabular}
\caption{Depth selects the emission. The count model's depth-weighting does work a
hard call cannot only in the intermediate band; at saturation the count model
\emph{is} a hard genotype.}
\label{tab:depth-regime}
\end{table}

%======================================================================
\section{The unified duration-aware three-state HMM engine}
\label{sec:engine}
%======================================================================
% Grounded in nilhmm/R/{engine,emissions,duration}.R. The compact version of this
% is the main-text Methods; here is the full specification.

Text~S2 established that ancestry is a spatially-correlated latent state that must
be inferred, not read off per site, and that the emission is the one axis that
should change with the data modality. The engine realizes exactly that: a single
duration-aware three-state HMM in which the emission and the duration/transition
prior are independent, swappable modules and everything else is shared. It is
data-agnostic --- $(\text{data},\text{params})\rightarrow\text{calls}$ --- so a
caller is nothing but a choice of (emission $\times$ duration $+$ priors)
(Text~S4). This is the local-ancestry HMM lineage
\citep{falush_2003,hapmix_2009,rfmix_2013,ancestryhmm_2017} specialized to a
biparental NIL design and decoded by standard forward--backward/Viterbi machinery
\citep{rabiner_1989}.

\subsection{States, observations, and the generative NIL model}
\label{sec:engine-model}

The hidden state at marker $t$ is the local ancestry
$s_t\in S=\{\REF,\HET,\ALT\}=\{0,1,2\}$ (recurrent-hom, heterozygous, donor-hom).
The observation is either a read-count pair $o_t=(n_t,a_t)$ (total depth $n_t$,
alternate/donor count $a_t$) or a hard genotype call $g_t\in\{0,1,2,\text{missing}\}$,
depending on the emission (Text~S2.4). A NIL genome is a piecewise-constant ancestry
track --- long donor blocks on a recurrent background --- so the modelling task is
to segment the marker sequence into runs of constant $s$.

\subsection{Emission axis}
\label{sec:engine-emission}

\paragraph{Count (BetaBinomial).} For read counts, the state-conditional emission is
a BetaBinomial on the alternate count, a direct parameterization of the per-site
genotype-likelihood sufficient statistic $(n_t,a_t)$ of Eq.~\eqref{eq:accumulate}
with overdispersion for mapping/paralog noise (the emission of the RTIGER
recombinant-genome genotyper, \citealp{rtiger_2023}),
\begin{equation}
\psi_s(o_t)=\BetaBin(a_t;\,n_t,\theta_s,\kappa)
=\binom{n_t}{a_t}\,
\frac{\Bfun\!\big(a_t+\kappa\theta_s,\;n_t-a_t+\kappa(1-\theta_s)\big)}{\Bfun\!\big(\kappa\theta_s,\;\kappa(1-\theta_s)\big)},
\label{eq:betabin}
\end{equation}
where $\theta_s$ is the state's expected alt fraction and $\kappa$ the concentration
(overdispersion; $\kappa\to\infty$ recovers the Binomial). Depth-zero markers
($n_t=0$) emit a flat likelihood. The default means are the textbook
$\theta=(\varepsilon,\,\tfrac12,\,1-\varepsilon)$, but under reference bias or the
outbred-donor dilution of Eq.~\eqref{eq:compressed-means} the true means are
compressed toward $\{0,\tfrac{d}{2},d\}$; the means are therefore optionally
EM-fittable (Text~S3.5), which is required for RNA/BRB-seq and recommended whenever
the donor-dose $d$ departs from $1$.

\paragraph{Categorical (genotype confusion).} At saturated depth a hard call
$g_t$ is sufficient (Text~S2.4), and the emission is a $3\times4$ categorical
confusion matrix over $\{0,1,2,\text{missing}\}$ (Holland's nNIL genotype model),
$\psi_s(g_t)=\Phi_{s,g_t}$, parameterized by the homozygote error $germ$, the
heterozygote error $gert$, the fraction $p$ of homozygous errors miscalled as het,
the missing rate $mr$, and the non-informative-marker rate $nir$; e.g.\ the
recurrent-hom row is
\begin{equation}
\Phi_{\REF,\cdot}=\big((1-germ)(1-mr),\;p\,germ\,(1-mr),\;(1-p)\,germ\,(1-mr),\;mr\big),
\label{eq:gtmat}
\end{equation}
with the het and donor-hom rows defined analogously (full matrix in the package
source). Raising $germ$/$gert$ toward the platform error rate absorbs isolated
miscalled markers as errors rather than opening spurious one-marker segments.

\subsection{Duration / transition axis}
\label{sec:engine-duration}

\paragraph{Geometric (memoryless).} The default transition has a per-marker switch
rate $r$: self-transitions $1-r$, and off-diagonal mass split between the two other
states in proportion to their prior frequencies (Text~S3.4),
\begin{equation}
A=\begin{pmatrix}
1-r & r\,\dfrac{f_1}{f_1+f_2} & r\,\dfrac{f_2}{f_1+f_2}\\[8pt]
r\,\dfrac{f_0}{f_0+f_2} & 1-r & r\,\dfrac{f_2}{f_0+f_2}\\[8pt]
r\,\dfrac{f_0}{f_0+f_1} & r\,\dfrac{f_1}{f_0+f_1} & 1-r
\end{pmatrix},
\qquad
\pi=(f_0,f_1,f_2),
\label{eq:transition}
\end{equation}
so the run-length in any state is geometric with mean $1/r$. The generation enters
the transition \emph{only} through the priors $f_1,f_2$ --- never as a baked-in
duration --- keeping the duration a genuine free axis.

\paragraph{Rigidity (Erlang state-expansion).} To enforce a hard minimum run length
$r$ (the RTIGER-style regularizer, \citealp{rtiger_2023,tiger_2015}), each
macro-state $m$ is expanded into a chain of
$r$ sub-states $S'=S\times\{1,\dots,r\}$: sub-states $1,\dots,r-1$ advance
deterministically (forcing the minimum run), and the terminal sub-state $r$ is free
--- it either self-loops to stay in $m$ (mass $A_{mm}$) or exits to sub-state $1$ of
another macro-state $m'$ (mass $A_{mm'}$),
\begin{equation}
T'_{(m,k)(m',k')}=
\begin{cases}
1 & m'=m,\ k'=k+1,\ k<r \quad\text{(forced advance)}\\
A_{mm} & m'=m,\ k=k'=r \quad\text{(free self-loop)}\\
A_{mm'} & m'\ne m,\ k=r,\ k'=1 \quad\text{(free exit)}\\
0 & \text{otherwise.}
\end{cases}
\label{eq:rigidity}
\end{equation}
All $r$ sub-states of a macro-state share its emission. This is a phase-type
(Erlang) sojourn; $r=1$ collapses Eq.~\eqref{eq:rigidity} back to
Eq.~\eqref{eq:transition}. It is a duration prior with a \emph{fixed}
priors-derived transition --- distinct from the full rigid-HMM estimator of the
RTIGER port (Text~S8), which additionally EM-learns $A$ and $\Psi$.

\paragraph{Explicit-duration (HSMM).} A genuine explicit-sojourn layer is reserved
but deliberately not used to bake in the design's fitted sojourn distribution, which
would make the duration circular with the breeding design it is meant to test.

\subsection{Design priors}
\label{sec:engine-priors}

The priors $(f_0,f_1,f_2)$ are the single-locus genotype frequencies of the
breeding design, obtained by iterating the backcross and selfing transition
matrices from the F$_1$ (Text~S5). They enter only through
Eqs.~\eqref{eq:transition}--\eqref{eq:rigidity} (as the transition's off-diagonal
weights and the start distribution $\pi$). For the two BZea designs,
$(f_0,f_1,f_2)=(0.844,0.0625,0.0938)$ for BC2S2 and $(0.859,0.0312,0.1094)$ for
BC2S3; both share the donor dosage $p=f_2+\tfrac12 f_1=0.125=2^{-(n_{bc}+1)}$,
invariant across selfing.

\subsection{Inference: Viterbi and EM for the emission means}
\label{sec:engine-inference}

Decoding returns the maximum-a-posteriori state path by Viterbi over
Eqs.~\eqref{eq:betabin}/\eqref{eq:gtmat} and
\eqref{eq:transition}/\eqref{eq:rigidity},
$\hat s_{1:T}=\argmax_{s_{1:T}}\ \pi_{s_1}\prod_{t}A_{s_{t-1}s_t}\,\psi_{s_t}(o_t)$,
with the rigidity sub-states mapped back to macro-states. When the emission means
are fittable, they are estimated by Baum--Welch EM with the transition held fixed:
the E-step computes the forward--backward posteriors $\gamma_s^t$, and the M-step is
the posterior-weighted alt fraction
\begin{equation}
\theta_s \;=\; \frac{\sum_t \gamma_s^t\,a_t}{\sum_t \gamma_s^t\,n_t},
\label{eq:em-means}
\end{equation}
pooled across a sample's chromosomes so the donor-state mean is estimated from all
donor evidence. Equation~\eqref{eq:em-means} is exact for a Binomial and is the
fixed-$\kappa$ mean update for the BetaBinomial; a state that is never visited keeps
its current mean. Only the means are learned --- the run-length parameter $r$ is set
by simulation calibration (Text~S6), not by likelihood, for the reasons given
there.

%======================================================================
\section{The caller family}
\label{sec:callers}
%======================================================================

The named callers are not separate programs but points in the engine's
$(\text{emission}\times\text{duration}+\text{priors})$ space
(Table~\ref{tab:callers}). Four ride the shared engine of Text~S3; two are faithful
ports of external biparental imputers, carried as baselines; and one is a
deliberately unsmoothed control. These biparental-population callers are the
low-coverage analogues of the population-scale local-ancestry and haplotype-imputation
HMMs \citep{falush_2003,hapmix_2009,rfmix_2013,ancestryhmm_2017,mosaic_2019,mach_2010,beagle_2007,stitch_2016,glimpse_2021},
whose common premise --- that linkage structure reduces per-marker genotype
uncertainty and thereby raises association power \citep{marchini_howie_2010} --- is
the same one Text~S7 formalizes.
\begin{table}[t]
\centering
\small
\begin{tabular}{@{}l l l l >{\raggedright\arraybackslash}p{4.6cm}@{}}
\toprule
caller & emission & duration & key param & role \\
\midrule
nNIL      & count / BetaBin      & geometric            & \texttt{rrate}      & count baseline; the self-transition \emph{is} the smoother \\
RTIGER    & count / BetaBin      & rigidity             & \texttt{rigidity}   & min.\ run length; robust at low depth (Julia-free port) \\
binhmm    & Gaussian / $1$\,Mb bins & sticky            & bin width           & pre-bin $\to$ 3-state Gaussian HMM; full SNP density (skim) \\
ATLAS     & categorical GT       & geometric            & call thresholds     & competitive-alignment RNA reads $\to$ hard GT (BRB/ASE) \\
\midrule
LB-Impute & external             & linkage-decay dist.\ & \texttt{recombdist} & Fragoso \emph{et al.}\ biparental imputer (ported baseline) \\
FSFHap    & external             & family-pooled EM     & \texttt{phet}       & TASSEL parent-call $+$ 5-state EM (ported baseline) \\
control   & per-site GL, no HMM  & ---                  & ---                 & het-prone no-HMM baseline; the het-excess reference \\
\bottomrule
\end{tabular}
\caption{The caller family as coordinates in the engine's two axes (top: shared
engine) plus external ports and a control (bottom). RTIGER and TIGER
\citep{rtiger_2023,tiger_2015}; LB-Impute \citep{lbimpute_2016}.}
\label{tab:callers}
\end{table}

The four shared-engine callers differ only along the two axes of Text~S3. \textbf{nNIL}
and \textbf{RTIGER} share the count emission and the same read-covered marker
support, differing only in duration (geometric vs.\ rigidity); \textbf{ATLAS} swaps
the emission to categorical because competitive-alignment RNA reads carry a hard
call rather than an unbiased allele count (expression and allele-specific expression
confound the count model); \textbf{binhmm} pre-bins to $\sim$1\,Mb before a 3-state
Gaussian HMM, the only track that keeps the full marker density on skim data.
\textbf{LB-Impute} and \textbf{FSFHap} are ported verbatim as external baselines ---
each carries its own transition/imputation model (a distance-dependent
recombination decay, and per-family EM haplotype imputation, respectively) but emits
the common segment schema, so they are directly comparable. The \textbf{control} is
the per-site genotype-likelihood MAP of Text~S2 with \emph{no} transition model; it
isolates exactly what the duration prior buys, and (being het-prone at low coverage,
Eq.~\eqref{eq:hwe-ratio}) is scored by its donor-called fraction on a known-recurrent
control region --- the ``het-excess'' reference the HMM callers must beat.

%======================================================================
\section{The \texttt{simcross} model and single-locus validation}
\label{sec:simcross}
%======================================================================

Calibration (Text~S6) needs a source of \emph{known} ancestry mosaics. We
forward-simulate the breeding design with \texttt{simcross}: the BZea BC2S3 (and the
bulked BC2S2 for skim) crossing scheme on the native B73 v5 consensus genetic map,
with crossover interference (Stahl model, $m=10$, $p=0$), then degrade each true
mosaic to a target modality by drawing read counts under the missingness law of
Text~S1 (per-marker coverage with floor $\pi$ and decay $k$) plus the per-read error
of Text~S2. Because the degradation reuses the empirically-fit $\miss(\cov)$ and its
scale-invariance (Text~S1.4), the simulated cohorts carry a realistic
coverage/missingness structure while their ancestry is exact truth.

Before trusting any caller's recovery on this truth, we verify that the simulator
itself is unbiased at the single-locus margin --- a circularity check on the truth
generator. The expected genotype frequencies follow from iterating the backcross and
selfing single-locus transition matrices from the F$_1$; for BC2S3,
\begin{equation}
\begin{aligned}
(f_{\REF},f_{\HET},f_{\ALT}) &= (0.859375,\ 0.03125,\ 0.109375),\\
\text{donor dosage } d &= f_{\ALT}+\tfrac12 f_{\HET} = 0.125 = 2^{-(n_{bc}+1)},
\end{aligned}
\label{eq:bc2s3-freqs}
\end{equation}
the dosage being invariant across selfing. We test whether the per-NIL genome
fractions (segment-length-weighted, in Mb) reproduce Eq.~\eqref{eq:bc2s3-freqs} with
a \emph{design-correct} statistic: a Hotelling $T^2$ on the $(\HET,\ALT)$ fractions
with the \emph{between-sample} empirical covariance,
\begin{equation}
T^2 = N\,(\hat{\mathbf f}-\mathbf f_0)^{\!\top} S_2^{-1}\,(\hat{\mathbf f}-\mathbf f_0)
\ \sim\ \chi^2_2,
\qquad
\mathrm{SE}(f_i)=\mathrm{sd}(f_i)/\sqrt{N},
\label{eq:hotelling}
\end{equation}
where the sample --- not the marker --- is the unit of replication. A na\"ive
per-marker binomial test would falsely reject a correct sampler, because linkage
makes the markers within a NIL massively over-dispersed relative to independence; the
effective information is only the $\sim\!33$ crossovers per genome, so the honest
standard error is $\mathrm{sd}(f_i)/\sqrt{N}$ over NILs. The simulator \emph{passes}:
all three state frequencies fall inside their $95\%$ Wald intervals and the donor
dosage matches $0.125$, so the truth generator is unbiased where the breeding design
makes a sharp prediction. The design-determined (non-Hardy--Weinberg) nature of this
variance is the standard logic of controlled-cross mapping populations
\citep{lander_botstein_1989,yu_2008,buckler_2009}, here for the TeoNAM-adjacent BZea
material \citep{chen2019teonam}; the crossover count and the resulting tract
structure follow the theory of junctions \citep{fisher_1954,stam_1980}.

A separate marker-resolution sweep isolates a \emph{truth-measurement} floor distinct
from caller error: measuring each NIL's fractions on progressively coarser marker
grids, the heterozygous fraction is undercounted (Cohen's
$d=(\bar f-f_0)/\mathrm{sd}(f)<0$) at wide spacing and relaxes into the negligibility
band $|d|<0.05$ once the grid is denser than $\sim$0.5\,Mb --- the resolution at
which segment endpoints, not the caller, limit accuracy.

%======================================================================
\section{Two-stage simulation calibration}
\label{sec:calibration}
%======================================================================

Each caller carries one free segmentation parameter --- nNIL's \texttt{rrate},
RTIGER's \texttt{rigidity}, LB-Impute's \texttt{recombdist} --- and we set it from the
simulated truth of Text~S5 rather than by eye. The search is two-stage: a coarse
log-spaced sweep brackets the optimum, then a golden-section refinement
(ratio $(\sqrt5-1)/2$) on the continuous $\log_{10}$ parameter locates it, all from a
single fit whose decodes are fanned across the grid. Calibration is done on
simulated cohorts and the chosen value is scored on a \emph{held-out} truth set.

The objective is a distributional match to the true donor-fragment structure, not a
likelihood. A called donor block matches a truth block when they reciprocally overlap
$\ge 50\%$ of each other's length; from the matched sets,
\begin{equation}
\begin{gathered}
\text{recall}=\frac{\#\text{matched truth}}{\#\text{truth}},\qquad
\text{precision}=\frac{\#\text{matched called}}{\#\text{called}},\\[2pt]
\mathrm{Dice}=\frac{2\,\text{prec}\cdot\text{rec}}{\text{prec}+\text{rec}},\qquad
\mathrm{FDR}=1-\text{precision},
\end{gathered}
\label{eq:dice}
\end{equation}
supplemented by the fragment-size Kolmogorov--Smirnov distance
$D=\sup_x|F_{\text{called}}(x)-F_{\text{truth}}(x)|$ between the called and true
donor-block-size distributions. Dice is maximized, $D$ and FDR minimized. The metric
penalizes both missed donor and spurious one-marker blocks (which, spanning $0$\,bp,
can never reach $50\%$ reciprocal overlap), unlike marker recall, which saturates and
rewards over-fragmentation.

Crucially, $\texttt{rrate}$ and $\texttt{rigidity}$ are chosen as
\emph{minimum-distance} estimators, not maximum-likelihood estimates: we do not
maximize a likelihood of the parameter but select the value whose \emph{called}
fragment-size distribution best matches the simulated-truth distribution (smallest
$D$, or highest Dice / lowest FDR). This matters because each parameter acts as a
low-pass filter whose cutoff can be turned up without bound, so ``maximize
smoothness/coherence'' is the wrong objective --- the smoothest caller is not the most
accurate (Text~S7). One therefore compares callers each at its \emph{calibrated}
cutoff,
\begin{equation}
\max_{\text{caller}}\ \rho\big(d;\ \lambda_{\text{calibrated}}\big)
\qquad\text{not}\qquad
\max_{\lambda}\ \rho(d;\lambda),
\label{eq:calib-not-max}
\end{equation}
and the biologically-correct cutoff is fixed externally by the fragment scale: the
recombination rate and crossing scheme set the tract-length distribution
\citep{fisher_1954,gravel_2012,liang_nielsen_2014}, so the ideal filter passes
wavelengths down to the shortest \emph{real} introgression fragment and rejects
shorter ones as noise. That the optimum is intermediate --- under-smoothing makes
false switches, over-smoothing erases true breakpoints --- is the bias--variance
trade-off \citep{geman_1992}, and the residual attenuation from imperfect ancestry
recovery is regression dilution \citep{frost_thompson_2000}. On the BZea BC2S2
benchmark the refined values are $\texttt{rrate}^\star\approx1.1\times10^{-3}$ (nNIL)
and $\texttt{rigidity}^\star=2$ (RTIGER); at those settings the held-out benchmark
reaches marker-level Dice $\approx0.99$ and donor-fragment Dice $\approx0.97$ at
$\mathrm{FDR}\approx0.01$. For the coverage sweep, each caller is recalibrated per
coverage by minimum donor-fragment FDR, and RTIGER's rigidity is feasibility-capped
so every chromosome retains more than $2r$ covered markers.

%======================================================================
\section{GWAS-landscape geometry}
\label{sec:gwas-geometry}
%======================================================================

Calibration (Text~S6) asserts that a better ancestry reconstruction yields a cleaner
association scan. This section makes that precise by reading a Manhattan plot as a
noisy realization of a spatial stochastic process, whose second-order geometry can
be used to score ancestry-inference quality. Each ingredient below is a
domain-specific instance of an established result; the contribution is the synthesis
into a falsifiable, three-axis caller diagnostic.

\subsection{Manhattan height as information}
\label{sec:noncentrality}

The plotted height $S(x)=-\log_{10}P(x)$ is the self-information (surprisal)
$-\log P$ of Shannon \citep{shannon_1948}, and for the association statistic
$Z=\hat\beta/\mathrm{SE}(\hat\beta)$ (with $Z^2\sim\chi^2_1$ under the null) the tail
gives $S(x)\approx Z^2/(2\ln 10)$. Its expectation under the alternative is the
\emph{non-centrality parameter} of the test \citep{sham_purcell_2014}, in the
mixed-model form used by plant and human GWAS \citep{yu_2006,zhou_stephens_2012},
\begin{equation}
\mathbb{E}[Z^2(x)] \;\propto\; N_{\text{eff}}(x)\,\frac{\mathrm{Var}\!\big(X(x)\big)\,\beta(x)^2}{\sigma^2},
\label{eq:ncp}
\end{equation}
i.e.\ effective sample size $\times$ variance of the segregating predictor $\times$
effect size $/$ residual variance. Peak height is thus information, and it scales
with the \emph{variance of the predictor} $X$ --- which in a mapping population is not
the Hardy--Weinberg value.

\subsection{Breeding-design variance versus Hardy--Weinberg}
\label{sec:designvar}

$\mathrm{Var}(X)=2p(1-p)$ holds only under random mating; a crossing scheme
correlates the within-individual alleles and inflates homozygosity, so the correct
variance is the genotype-frequency variance $\mathbb{E}[A^2]-\mathbb{E}[A]^2$ over the
design's realized frequencies \citep{lander_botstein_1989,yu_2008}. For the BC2S3
dosage $A\in\{0,1,2\}$ with Eq.~\eqref{eq:bc2s3-freqs} ($p=0.125$),
\begin{equation}
\mathbb{E}[A]=0.25,\quad \mathbb{E}[A^2]=0.46875,\quad
\mathrm{Var}(A)=0.40625 \approx 1.86\times \underbrace{2p(1-p)}_{0.21875},
\label{eq:vara}
\end{equation}
and collapsing to a two-state (donor vs.\ recurrent) contrast gives
$p_{\REF}p_{\ALT}=0.109$. The design nearly doubles the segregating variance relative
to the HWE value one would naively insert into Eq.~\eqref{eq:ncp} --- a fixed,
computable property of the pedigree, not of an allele-frequency equilibrium.

\subsection{Ancestry as predictor; low coverage as measurement error}
\label{sec:predictor}

The causal chain is genotype$\rightarrow$phenotype; local ancestry $A(x)$ is a
\emph{correlate} of genotype, so mapping on ancestry improves estimation of a
genotype-derived predictor rather than asserting ancestry is causal --- the logic of
admixture mapping \citep{mckeigue_1998,patterson_2004,winkler_2010}. The engine's
output is the HMM posterior expectation $\hat X=\mathbb{E}[X\mid\text{data},A]$
\citep{rabiner_1989,falush_2003}. Low coverage enters as measurement error on the
predictor, $\tilde X = X+\epsilon$; the HMM reduces $\mathrm{Var}(\hat X - X)$ below
$\mathrm{Var}(\tilde X - X)$, and the association effect is attenuated by the
reliability ratio
\begin{equation}
\hat\beta = \beta\,\frac{\mathrm{Cov}(\hat G,G)}{\mathrm{Var}(\hat G)},
\label{eq:attenuation}
\end{equation}
the classical regression-dilution correction \citep{frost_thompson_2000}, whose
attenuation vanishes as $\hat G\to G$. This is the mechanism by which better
imputation raises power \citep{marchini_howie_2010,mach_2010,beagle_2007,stitch_2016,glimpse_2021},
and it is what distinguishes ``the peak cleaned up because the predictor improved''
from ``the peak cleaned up because it was smoothed'' (Text~S7.7): attenuation acts on
the \emph{covariance}, not on spatial roughness.

\subsection{The peak as a spatial object}
\label{sec:spatial}

At a marker in incomplete LD ($r^2(d)$) with the causal variant, the statistic
attenuates as $\mathbb{E}[Z^2(d)]\propto N\,\mathrm{Var}(X)\beta^2 r^2(d)/\sigma^2$
--- the theory of indirect (tag-marker) association \citep{pritchard_2001} --- so a
point causal effect is spread into a peak by the correlation kernel, the spatial
statement of the interval-mapping LOD profile \citep{lander_botstein_1989}. A QTL
therefore produces \emph{spatially autocorrelated} evidence: with
$\rho(d)=\mathrm{cor}\big(S(x),S(x+d)\big)$ positive and decaying with distance ---
spatial autocorrelation \citep{moran_1950,geary_1954} and Tobler's first law
\citep{tobler_1970}. Summaries are the lag-one autocorrelation $\rho_1$ and the
integrated $Q=\sum_k\rho(k)$. Because the autocovariance and the power spectrum are a
Fourier-transform pair --- the Wiener--Khinchin theorem \citep{wiener_1930,khinchin_1934}
--- high-frequency spectral energy and $\rho_1$ measure the same second-order
structure and rank callers in reverse order; the band-limited power estimate follows
standard spectral analysis \citep{percival_walden_1993} and the wavelet detail-energy
measure of local roughness the multiresolution decomposition of \citet{mallat_1989}.

\subsection{Correlation length equals the tract scale}
\label{sec:tractscale}

The correlation length of a \emph{mapping-population} scan is set by recent
recombination --- the cross's own meioses --- not by historical population LD. LD as
genotype correlation and its generation by finite-population recombinational history
is \citet{hill_robertson_1968}; the per-meiosis decay $C(d)\approx(1-2r)$ and the
segment-homozygosity relation is \citet{sved_1971}. Hence
\begin{equation}
L_{\text{signal}} \approx L_{\text{ancestry tract}},
\label{eq:lengthscale}
\end{equation}
the GWAS-profile correlation length equals the inherited ancestry-tract scale, fixed
by the theory of junctions and the admixture-tract-length distribution
\citep{fisher_1954,stam_1980,gravel_2012,liang_nielsen_2014,pool_nielsen_2009}. The
shortest \emph{real} introgression fragment is then a Nyquist-style resolution floor:
the shortest wavelength a caller should preserve --- exactly the biological cutoff
that Text~S6 calibrates toward.

\subsection{The band-limited noise metric and caller ranking}
\label{sec:noisiness}

Operationally we score the roughness of each caller's $-\log_{10}P$ track within a
noise band of wavelengths $(R,P_{10})$: $R$ is the median adjacent-marker spacing of
the inference grid and $P_{10}$ the $10$th percentile of resolvable teosinte-tract
lengths (on the 118K panel, $R\approx77$\,kb, $P_{10}\approx181$\,kb). Three
equivalent lenses (Text~S7.4) are computed per chromosome: the FFT band-power
fraction $\texttt{fft\_band}$, the Haar wavelet band-detail fraction
$\texttt{wav\_band}$, and the lag-one autocorrelation $\texttt{acf\_lag1}$. As
Wiener--Khinchin predicts, $\texttt{fft\_band}$ (ascending $=$ cleaner) and
$\texttt{acf\_lag1}$ (descending $=$ more coherent) rank the callers in reverse order
(Table~\ref{tab:noisiness}): the duration-aware HMMs \texttt{binhmm}/\texttt{rtiger}
are least noisy and the no-HMM \texttt{control}/\texttt{nNIL} noisiest.
\begin{table}[t]
\centering
\begin{tabular}{l c c c}
\toprule
caller & \texttt{fft\_band} & \texttt{wav\_band} & \texttt{acf\_lag1} \\
\midrule
binhmm   & $0.0004$ & $0.0008$ & $0.999$  \\
rtiger   & $0.0006$ & $0.0012$ & $0.9988$ \\
lbimpute & $0.0020$ & $0.0065$ & $0.9932$ \\
control  & $0.0179$ & $0.0452$ & $0.9586$ \\
nnil     & $0.0179$ & $0.0445$ & $0.9597$ \\
\bottomrule
\end{tabular}
\caption{Spatial-coherence ranking of the STAM $-\log_{10}P$ track at
$\lambda=\infty$ on the 118K panel. $\texttt{fft\_band}$/$\texttt{wav\_band}$
(band-limited high-frequency energy, lower $=$ cleaner) and $\texttt{acf\_lag1}$
(lag-one autocorrelation, higher $=$ more coherent) rank in reverse order, as
Wiener--Khinchin requires.}
\label{tab:noisiness}
\end{table}

\subsection{Recovery versus smoothing: the H1/H2 falsification}
\label{sec:h1h2}

A cleaner scan alone is ambiguous, because any strong low-pass filter inflates
coherence (H1: model-imposed smoothing) --- the over-/under-fit optimum is the
bias--variance trade-off \citep{geman_1992}. The alternative (H2: authentic mosaic
recovery) predicts that reconstruction accuracy and mapping performance peak
\emph{together} at the true ancestry-block scale. We separate them with a three-axis
test over the caller parameter $\theta$: ancestry accuracy $E_A(\theta)=P(\hat A_\theta\ne A_{\text{true}})$,
landscape coherence $\rho(\theta)$ (Text~S7.6), and mapping performance $M(\theta)$.
H1 predicts $\argmax_\theta\rho(\theta)\ne\argmin_\theta E_A(\theta)$; H2 predicts
$\argmin_\theta E_A(\theta)\approx\argmax_\theta M(\theta)$. The central,
falsifiable claim tested by the coverage sweep is therefore: \emph{the caller
parameters that best reconstruct the Mendelian ancestry mosaic produce the most
accurate GWAS landscape} --- which is why calibration targets the fragment-scale
minimum-distance fit (Text~S6), not maximal smoothness.

%======================================================================
\section{Rigid-HMM derivations (the RTIGER port)}
\label{sec:rhmm}
%======================================================================

The \texttt{rtiger} caller is a Julia-free R/Rcpp port of the robust rigid-HMM (rHMM)
of \citet{rtiger_2023}, which itself extends the TIGER genotyping-by-sequencing
pipeline \citep{tiger_2015}. Unlike the fixed-transition Erlang duration layer of
Text~S3.3 (which regularizes run length but does not learn the transition), the rHMM
\emph{EM-fits} the transition $A$ and the emission $\Psi$ under a hard minimum-run
constraint, and selects the rigidity $r$ by an analytic minimum-distance criterion.
We reproduce the model and estimator in condensed form; the full proofs are in the
supplement of \citet{rtiger_2023}.

\subsection{The $r$-rigid HMM and its augmented-state equivalent}
\label{sec:rhmm-model}

A state path $\Ss=(s_1,\dots,s_T)$ is \emph{$r$-rigid} if every maximal run of
identical states has length $\ge r$; let $\Sr$ be the set of such paths. The rHMM is
a standard HMM restricted to $\Sr$,
\begin{equation}
P_{r\mathrm{HMM}}(\Oo,\Ss;\Theta_r):=P_{\mathrm{HMM}}(\Oo,\Ss\mid \Ss\in\Sr;\Theta)
=\prod_{t=1}^{T}\psi_{s_t}(o_t)\ \prod_{t=1}^{T}P\!\big(s_t\mid s_{t-1},\Ss\in\Sr\big),
\label{eq:rhmm}
\end{equation}
where the constrained transition is $a_{jk}$ once a run has reached length $r$ and is
forced to stay ($\delta_{jk}$) otherwise. The $r$-th-order chain is made first-order
by the same state expansion as Eq.~\eqref{eq:rigidity}: an augmented space
$S'=S\times\{1,\dots,r\}$ in which $(s,j)$ records a run of length $j$ in state $s$,
carrying a standard HMM $\Theta'_r$ with $P_{\mathrm{HMM}}(\Oo,\Ss;\Theta_r)=
P_{\mathrm{HMM}}(\Oo,\Ss'\mid j_T=r;\Theta'_r)$. This reduces inference from the
naive $|S|^r$ to $O(|S|^2 T)$.

\subsection{Rigid forward--backward (E-step)}
\label{sec:rhmm-estep}

The recursions run on $\Theta'_r$ and reuse the windowed emission product
\begin{equation}
\Psi^{t}_k \;=\!\!\prod_{u=\max(t-r+1,1)}^{\min(t,T)}\!\!\psi_k(o_u),
\label{eq:psiwindow}
\end{equation}
so that the terminal ($j=r$) forward mass telescopes,
\begin{equation}
\alpha'^{\,t}_{(v,r)}=\sum_{u\ne v}\alpha'^{\,t-r}_{(u,r)}\,a_{uv}\,\Psi^{t}_v
+\alpha'^{\,t-1}_{(v,r)}\,a_{vv}\,\psi_v(o_t),
\label{eq:fwd}
\end{equation}
with the analogous backward recursion. Collapsing the augmented posteriors back to
macro-states gives the per-marker $\gamma^t_k=P(s_t=k,\Oo)$ and the run-boundary
transition counts $\zeta^t_{jk}$ needed for the M-step (full definitions in
\citealp{rtiger_2023}).

\subsection{M-step updates}
\label{sec:rhmm-mstep}

Maximizing the EM lower bound under the simplex constraints gives closed forms for
the transition and initial distribution and a method-of-moments update for the
BetaBinomial emission,
\begin{equation}
a_{jk}=\frac{\sum_t \zeta^t_{jk}}{\sum_{k'}\sum_t \zeta^t_{jk'}},
\qquad
\pi_j=\gamma^1_j,
\qquad
m_s=\frac{a_s}{a_s+b_s}=\frac{\sum_t k_t\,\gamma^t_s}{\sum_t n_t\,\gamma^t_s},
\label{eq:mstep}
\end{equation}
where $(k_t,n_t)$ are the alternate and total counts and the emission's expected
alt-fraction $m_s$ is the posterior-weighted allele ratio --- the same estimator as
the engine's fixed-transition mean update Eq.~\eqref{eq:em-means}. The remaining
concentration $\tau_s$ (so $a_s=m_s\tau_s$, $b_s=(1-m_s)\tau_s$) is set by a 1-D line
search; multiple chains contribute additively to the pooled sufficient statistics.

\subsection{Minimum-distance rigidity tuning}
\label{sec:rhmm-rigidity}

The rigidity $r$ is chosen not by likelihood but by controlling segmentation error.
For a segment of length $u$ flanked by another state, the log-likelihood contrast
between the true and a competing state, $\Delta_{u,x,y}$, is sampled by Monte Carlo
under a Poisson-coverage emission $p_{k,n,z}=\Pois(n;c)\,\psi_z(k;n)$, giving the
false-positive and false-negative segment rates
\begin{equation}
\mathrm{FPR}_{x\to y}(u;\Theta_r)=P\!\big(\Delta_{u,x,y}+C_{u,x,y}>0\big),
\qquad
\mathrm{FNR}_{x\leftarrow y}(u;\Theta_r)=P\!\big(\tilde\Delta_{m,x,y}\le 0\big),
\label{eq:fprfnr}
\end{equation}
from which the expected numbers of false and missed segments per sample,
$\mathrm{SE}_+(\Theta_r)$ and $\mathrm{SE}_-(\Theta_r)$, follow by a constrained
optimization over segment-length compositions. The chosen $r$ minimizes the total
expected segmentation error --- a \emph{minimum-distance} selection on the segment
statistics, not an MLE of $r$. This is the analytic counterpart of the engine's
empirical route (Text~S6): where \citet{rtiger_2023} predict the segment error from
$\Theta_r$ and a coverage model, nilHMM's \texttt{calibrate\_r} measures the
fragment-size distance directly against \texttt{simcross} truth
(Eq.~\eqref{eq:dice}). Both estimate the run-length scale by matching the true
segment geometry, consistent with the bias--variance and tract-length arguments of
Text~S6--S7; neither treats $r$ as a likelihood-maximizing parameter.

%======================================================================
\section{Pedigree-coupled ancestry calling}
\label{sec:pedigree}
%======================================================================

Every caller of Text~S3--S4 decodes each individual's chromosome independently. In
a structured breeding population that leaves information on the table: siblings of a
selfing family inherit overlapping donor blocks from a shared ancestor, so their
genotypes are correlated through the pedigree, not just along the chromosome. The
pedigree caller adds that missing coupling --- a \emph{third} structural axis beyond
the emission and duration of Text~S3, namely inter-individual coupling through the
known family tree. It is the family-pooled analogue of the single-chain callers (in
the same spirit as the family-pooled FSFHap of Text~S4), and it reduces the
predictor error of Text~S7.3 using Mendelian structure the single-chain callers
ignore.

\subsection{The pedigree $\times$ genome grid and its two couplings}
\label{sec:ped-grid}

Nodes are individuals (genotyped leaves and latent ungenotyped ancestors); the
hidden state of node $v$ at marker $t$ is the donor dosage
$s^v_t\in\{\REF,\HET,\ALT\}=\{0,1,2\}$ as before. Two factors couple the grid.

\paragraph{Along the chromosome (per node).} A relax-to-stationary transition
builds ancestry blocks while relaxing to a node-specific stationary $\pi$,
\begin{equation}
A_{xy} = (1-r_{\mathrm{eff}})\,\delta_{xy} + r_{\mathrm{eff}}\,\pi_y,
\qquad
r_{\mathrm{eff}} = \tfrac12\big(1-(1-2r_t)^{m_v}\big),
\label{eq:ped-transition}
\end{equation}
where $r_t$ is the per-interval recombination fraction (Text~S3.4 / Haldane on the
map, or $\texttt{rrate}\cdot\Delta\mathrm{bp}$) and $m_v$ is the number of meioses
accumulated from the founder to node $v$; $r_{\mathrm{eff}}$ is the probability of
an odd number of crossovers over those $m_v$ meioses and saturates correctly at
$\tfrac12$. The stationary satisfies $\pi A=\pi$ for any $\pi$.

\paragraph{Across the pedigree (parent $\to$ child).} A single selfing meiosis
transmits dosage through the kernel
\begin{equation}
T^{\mathrm{sib}}=
\begin{pmatrix}
1 & 0 & 0\\
\tfrac14 & \tfrac12 & \tfrac14\\
0 & 0 & 1
\end{pmatrix}
\qquad (\text{row}=\text{parent dosage},\ \text{col}=\text{child}),
\label{eq:tsib}
\end{equation}
so a homozygous locus stays fixed in all descendants and heterozygosity halves each
selfing generation --- exactly the segregation the pedigree information rides on.

\paragraph{Where the genotype-frequency prior enters.} Only at the founder. With the
donor genome fraction $q=f_2+\tfrac12 f_1$ (Text~S3.4; selfing-invariant), the
founder is heterozygous for all its donor content, giving the founder prior and
chain stationary
\begin{equation}
\pi_0=\big(1-2q,\ 2q,\ 0\big)
\qquad(\text{BC2}\Rightarrow(0.75,\,0.25,\,0)),
\label{eq:founder-prior}
\end{equation}
while every non-root node takes a \emph{uniform} stationary and marker-0 prior. The
deeper generation marginals $\pi_t\approx(0.859,0.031,0.109)$ are \emph{emergent},
not injected: applying $T^{\mathrm{sib}}$ to $\pi_0$ yields $\pi_1,\pi_2,\dots$, so
the transmission messages carry them downward. Re-injecting $\pi_t$ as each node's
chain stationary would double-count the prior against the transmission message and
over-suppress the conditionally-correct heterozygosity of a selfing child
($(0.25,0.5,0.25)$ off a HET parent) toward the $3\%$ population marginal --- so the
non-root chains are deliberately marginal-neutral.

\subsection{Structured loopy belief propagation}
\label{sec:ped-bp}

The grid is loopy --- adjacent markers across a parent--child edge form a 4-cycle
--- so exact inference is intractable and we use structured loopy belief propagation
\citep{pearl_1988,kschischang_2001,yedidia_2005}. Each node's chromosome chain is
solved \emph{exactly} by a tilted forward--backward \citep{rabiner_1989}: the
effective per-marker likelihood is the node's emission times a \emph{tilt} equal to
the product of all incoming pedigree messages,
$\tilde\psi^v_t(x)=\psi^v_t(x)\prod_{u\in\mathcal N(v)}\mu_{u\to v}(t,x)$; latent
ancestors set $\psi\equiv1$. The pedigree edges are the loopy part, passed by damped
cavity messages. A message contracts $T^{\mathrm{sib}}$ over the sender's cavity
belief (the belief with the recipient's own message divided out),
\begin{equation}
\mu_{v\to p}(t,x_p)\propto\sum_{x_v} T^{\mathrm{sib}}_{x_p x_v}\,\mathrm{cav}_{v\setminus p}(t,x_v),
\qquad
\mu_{p\to v}(t,x_v)\propto\sum_{x_p} T^{\mathrm{sib}}_{x_p x_v}\,\mathrm{cav}_{p\setminus v}(t,x_p),
\label{eq:messages}
\end{equation}
(up = child$\to$parent, down = parent$\to$child), updated with a damped geometric
step $\mu\leftarrow\mu^{\,1-\lambda}\,\mu_{\mathrm{new}}^{\,\lambda}$ ($\lambda\approx0.5$)
over post-order-then-pre-order sweeps until the maximum message change falls below
tolerance. Decoding is \emph{marginal-MAP}: the per-marker state is
$\hat s^v_t=\argmax_x \gamma^v_t(x)$ of the converged posterior belief --- a
posterior-marginal argmax, not the Viterbi joint path of the single-chain callers.

\subsection{Emission and the two front doors}
\label{sec:ped-emission}

The emission is the same swappable axis as the rest of the engine (Text~S3.2),
which is what lets one kernel serve both a de novo caller and a refinement:
\begin{itemize}
\item \textbf{Read counts} $\to$ the BetaBinomial \texttt{emission\_count}
(Eq.~\eqref{eq:betabin}); a zero-coverage marker has $n_t=0$ and emits flat, so the
pedigree fills it in weighted by how confident the relatives are. This is the
depth-aware \emph{de novo} caller, \texttt{call\_ancestry(caller="pedigree")}.
\item \textbf{Hard calls} (a per-marker \texttt{state}/$g$ column) $\to$ the
categorical \texttt{emission\_gt} (Eq.~\eqref{eq:gtmat}); missing is the explicit
$g=3$ column, again flat. This is the refinement path, wrapped as
\texttt{refine\_ancestry()}.
\end{itemize}
In both channels ``no data'' collapses to a flat per-marker likelihood --- $n_t=0$
under the count emission, the $g=3$ column under the categorical --- which is
precisely what lets the pedigree and the chain supply the missing state (Text~S2.4).
The hard-call channel discards per-marker depth (a call is a call, Text~S2), so the
count channel is preferred whenever read counts are available; the categorical
channel is for correcting an existing mosaic or a VCF of hard genotypes.

\subsection{Complexity, approximations, and expected gains}
\label{sec:ped-complexity}

Each sweep costs $O(|V|\,M\,K^2)$ with $K=3$, so the fit is linear in both pedigree
size and marker count and converges in a few damped sweeps; families are fully
separable and embarrassingly parallel. Two approximations are stated honestly. First,
using $T^{\mathrm{sib}}$ as a \emph{per-marker} factor drops \emph{transmission
linkage} --- a heterozygous parent transmits a donor segment as a coherent block, not
marker-by-marker, so the cross-sibling block-coherent depth pooling is discarded and
only marker-by-marker marginal pooling survives (the phased copy-switch kernel is the
fix, deferred). Second, the per-site cavity is exact at its own site but a message's
influence propagated through the chain to neighboring sites is not removed, so a
neighbor's evidence leaks around the 4-cycle; damping controls this in practice
\citep{yedidia_2005}. Because the coupling reduces the ancestry-predictor error
(Text~S7.3; regression-dilution attenuation, \citealp{frost_thompson_2000}) using the
IBD block structure siblings share by descent \citep{fisher_1954,stam_1980}, the
gains are largest exactly where the single-chain callers are weakest: at low
coverage, near breakpoints, and for leaves with many genotyped siblings.

\bibliographystyle{plainnat}
\bibliography{references}

\end{document}